\manualchapter{Control reference}{sec:controls}
Panel map: \textbf{1} envelope sliders, Loop and Rate Scale;
\textbf{2} tuning and modulation knobs;
\textbf{3} envelope CV inputs;
\textbf{4} performance inputs and audio output.

The six top sliders and their CV sockets shape the envelope. The lower knobs
set tuning, feedback and modulation; the bottom sockets handle performance
signals and audio.

\subsection{Envelope sliders, left to right}
\begin{tabularx}{\textwidth}{@{}l l X@{}}
\toprule Control & Range; default & Result \\\midrule
AR & 1--31; 31 & Attack rate; larger values rise faster. \\
TL & 0--100; 100 & Peak level; larger values are louder. \\
D1R & 0--31; 0 & First decay rate; zero holds the peak. \\
SL / T1L & 0--15; 15 & Level at which the second decay begins. \\
D2R & 0--31; 0 & Second decay rate; zero holds this level. \\
RR & 0--15; 15 & Release rate after gate falls. \\
\bottomrule
\end{tabularx}

These are chip rate and attenuation settings, not seconds or linear amplitude
percentages. Rate scaling (0--3, default 0) makes envelopes respond more quickly
at higher pitches. Loop is off by default; it enables the existing simplified repeating
SSG envelope contour. Use it for repeating contours after first establishing a gated sound.
The hardware envelope terminology is retained in the reference figure and
in \citet{genesis-sound-software-manual}.

Each envelope CV adds $mV/8$ to its slider, where $m$ is that slider's maximum.
The sum is truncated to an integer and clamped to the control range. For
example, with AR at 15, $+4$\,V adds 15.5 before truncation, giving 30.
This is an offset; a 0\,V cable leaves the slider's setting intact.
